% !TEX program = platex
% 数学ノート集・共通テンプレート v2（積分論2026年度第1部の版面を参考）
% A5判。章内で定義・定理・例・注意等の番号を共通化。
% コンパイル： platex -interaction=nonstopmode math_note_template_v2.tex (2回)
%              dvipdfmx math_note_template_v2.dvi
\documentclass[a5paper,10pt,onecolumn]{jsarticle}
\usepackage[dvipdfmx,top=22mm,bottom=20mm,inner=19mm,outer=18mm,headheight=13pt,headsep=6mm]{geometry}
\usepackage{amsmath,amssymb,amsthm,mathtools}
\usepackage{enumitem,fancyhdr,titlesec}
\usepackage[dvipdfmx,bookmarks=true,colorlinks=false]{hyperref}
\usepackage{pxjahyper}

% ----- 差し替える文書情報 -----
\newcommand{\NoteTitle}{計算線形代数}
\newcommand{\NoteSeries}{数学ノート集}

% ----- 文字・組版 -----
% 参考PDFの質素な明朝＋Computer Modernの印象を引き継ぎながら，
% A5判で窮屈にならないよう本文の行間を少し確保する．
\linespread{1.13}
\setlength{\parindent}{1zw}
\setlength{\parskip}{0pt}
\setlength{\abovedisplayskip}{6pt plus 2pt minus 2pt}
\setlength{\belowdisplayskip}{6pt plus 2pt minus 2pt}
\setlength{\abovedisplayshortskip}{3pt plus 2pt minus 1pt}
\setlength{\belowdisplayshortskip}{4pt plus 2pt minus 1pt}
\setlist[enumerate]{topsep=3pt,itemsep=1pt,parsep=0pt,leftmargin=2.4zw}
\titleformat{\section}{\normalfont\normalsize\bfseries}{\thesection}{0.8zw}{}
\titlespacing*{\section}{0pt}{2.0ex plus .4ex}{.9ex}
\titleformat{\subsection}{\normalfont\normalsize\bfseries}{\thesubsection}{0.7zw}{}
\titlespacing*{\subsection}{0pt}{1.5ex plus .3ex}{.6ex}
\raggedbottom

% ----- 柱とページ番号：ページ上部、罫線なし -----
\fancypagestyle{mathnote}{%
  \fancyhf{}%
  \fancyhead[L]{\footnotesize \NoteSeries\ （\NoteTitle）}%
  \fancyhead[R]{\footnotesize \thepage}%
  \renewcommand{\headrulewidth}{0pt}%
  \renewcommand{\footrulewidth}{0pt}%
}
\pagestyle{mathnote}

% ----- 節内通し番号の定義・定理群 -----
\newtheoremstyle{mathnotetheorem}
 {4.5pt}{4.5pt}{\normalfont}{}%
 {\bfseries}{}{0.65zw}%
 {\thmname{#1}\thmnumber{ #2}\thmnote{（#3）}}
\theoremstyle{mathnotetheorem}
\newtheorem{theorem}{定理}[section]
\newtheorem{definition}[theorem]{定義}
\newtheorem{lemma}[theorem]{補題}
\newtheorem{proposition}[theorem]{命題}
\newtheorem{corollary}[theorem]{系}
\newtheorem{example}[theorem]{例}
\newtheorem{remark}[theorem]{注意}
\renewcommand{\proofname}{証明}
\renewcommand{\qedsymbol}{$\square$}
\makeatletter
\renewenvironment{proof}[1][\proofname]{%
  \par\pushQED{\qed}\normalfont\topsep4\p@\@plus2\p@\relax
  \trivlist\item[\hskip\labelsep\bfseries #1]\ignorespaces
}{\popQED\endtrivlist\@endpefalse}
\makeatother

% ----- 使うノートで適宜拡張する記号 -----
\newcommand{\R}{\mathbb{R}}
\newcommand{\norm}[1]{\lVert #1\rVert}
\newcommand{\abs}[1]{\lvert #1\rvert}

\begin{document}
\section*{到達目標}
整数行列を可逆な整数基本操作でSmith標準形に直し，
整数線形方程式の可解性と商加群を計算できるようになる．
前提は整数のEuclid算法と行列の基本操作である．

\section{整数行列の分数なし消去}
$A\in\mathbb Z^{n\times n}$ を考える．
分数の発生を避けるためBareissの更新を使う．
\begin{definition}
$A^{(0)}=A$，$d_{-1}=1$ と置く．
行と列の交換を適宜許し，$k$ 段目の非零pivotを
$d_k=a^{(k)}_{kk}$ としたとき，$i,j>k$ に対して
\[
 a^{(k+1)}_{ij}=
 \frac{d_k a^{(k)}_{ij}-a^{(k)}_{ik}a^{(k)}_{kj}}{d_{k-1}}
\]
と更新する．
\end{definition}
\begin{theorem}[Bareiss更新の整除性]
先頭の主小行列式が各段階で零でないとする．
このとき更新式の分子は $d_{k-1}$ で割り切れ，
途中の成分は整数である．最後のpivotは，実施した行・列の交換総数を $s$ とすると $(-1)^s\det A$ に等しい．
\end{theorem}
\begin{proof}
$k$ 段階での行列要素は，もとの行列の先頭 $k$ 行・列と
行 $i$，列 $j$ の交差小行列式と同一視できることを帰納的に示す．
$k=0$ では元の成分である．
次の段階の $2\times2$ 型の行列式恒等式
\[
 \det(B)\det(B_{\widehat{i},\widehat{j}})
 =\det(B_{\widehat i,\widehat i})
  \det(B_{\widehat j,\widehat j})
 -\det(B_{\widehat i,\widehat j})
  \det(B_{\widehat j,\widehat i})
\]
（Desnanot--Jacobi恒等式）を適用すると，
更新式の分子が前段階pivotと次の交差小行列式の積になる．
この恒等式は可逆な左上ブロックについてSchur補行列の
$2\times2$ 行列式を計算して得られ，多項式恒等式だから一般にも成り立つ．
よって整除性と不変量が帰納的に従う．行・列の交換がなければ最後は全行列式に等しい．
交換総数を $s$ とした場合，行列式の符号は $(-1)^s$ 倍になる．
\end{proof}
\section{Smith標準形の基本不変量}
\begin{proposition}
整数行列 $A\in\mathbb Z^{m\times n}$ が
可逆な整数行列 $P,Q$ を用いて
\[
 PAQ=\operatorname{diag}(d_1,\ldots,d_r,0,\ldots)
\]
と変形され，$d_i>0$，$d_i\mid d_{i+1}$ とする．
$\Delta_k$ を $A$ の全 $k$ 次小行列式の最大公約数とすると
\[\Delta_k=d_1\cdots d_k\quad(1\le k\le r).\]
\end{proposition}
\begin{proof}
Cauchy--Binet公式によって $PAQ$ の各 $k$ 次小行列式は
$A$ の $k$ 次小行列式の整数係数線形結合である．
逆変換にも同じ議論を適用できるため，小行列式の生成する整数イデアルは不変である．
対角形では全 $k$ 次小行列式が $d_1\cdots d_k$ の倍数であり，
先頭 $k$ 次小行列式がその値そのものなので最大公約数は表示式に等しい．
\end{proof}
\begin{remark}
一般のSmith標準形の存在算法は既存の15.010で扱われている．
ここでは不変因子の計算法と厳密な行列式計算を主題とした．
\end{remark}


\section{Smith標準形の構成算法}
\begin{theorem}[Euclid整域上のSmith算法]
$R$ をEuclid整域とする．任意の $A\in R^{m\times n}$ に対し
可逆行列 $P\in\mathrm{GL}_m(R)$，$Q\in\mathrm{GL}_n(R)$ が存在して
\[
 PAQ=\operatorname{diag}(d_1,\ldots,d_r,0,\ldots),
 \qquad d_i\mid d_{i+1}
\]
となる．しかもEuclid除法・行列の基本行列操作を用いて
有限回でこれを構成できる．
\end{theorem}
\begin{proof}
零行列の場合は明らかである．非零成分を行列の左上に置き，
その値を $a\ne0$ とする．
第一列に $a$ で割り切れない成分 $b$ があれば，
その行の $b$ を $a$ で割って $b=qa+r$，$0\ne r$ と書き，
行から $q$ 倍の第一行を引き，次に行交換して $r$ をpivotにする．
Euclid関数は厳密に減少する．第一行についても列操作で同様に行う．
これを繰り返すと左上 $a$ は第一行と第一列のすべての成分を割る．
行と列の基本操作でそれらを零にできる．
もし右下部分行列に $a$ で割り切れない成分 $b$ があれば，
その行を第一行に加えて第一行に $b$ を現し，
先ほどと同じ除算操作によって，左上pivotをより小さいEuclid関数の
非零元へ取り替えられる．
Euclid関数は非負整数値なので，この減少は有限回で停止する．
停止時の左上pivot $a$ は右下部分行列の全成分を割る．
右下部分行列について帰納的に同じ手続きを施すと，
その対角成分も $a$ の倍数のまま保たれる．
従って対角化と $d_i\mid d_{i+1}$ が得られる．
全操作は基本行・列操作であり，それぞれ可逆であるから
対応する $P,Q$ が存在する．
\end{proof}
\begin{corollary}[整数行列の不変因子の復元]
整数行列 $A$ の階数を $r$ とし，$\Delta_k$ を
$k$ 次小行列式全体の最大公約数，$\Delta_0=1$ とする．
このときSmith対角成分は符号を除いて
$ d_k=\Delta_k/\Delta_{k-1}$（$1\le k\le r$）で一意に定まる．
\end{corollary}
\begin{proof}
前定理で $R=\mathbb Z$ と取ればSmith形が存在する．
先に証明した小行列式イデアルの不変性より
$\Delta_k=|d_1\cdots d_k|$ が成り立つ．
隣り合う式の商を取れば表示式になり，
従って不変因子の絶対値は一意に定まる．
\end{proof}

\section{Smith不変因子のビット長の上界}
\begin{theorem}[最大小行列式による上界]
$A\in\mathbb Z^{m\times n}$ の階数を $r\ge1$ とし，
全成分の絶対値が $B\ge1$ 以下であるとする．
Smith標準形の非零不変因子を正整数
$d_1\mid\cdots\mid d_r$ とすると
\[
 d_1\cdots d_r\le r^{r/2}B^r,
 \qquad \log_2d_i\le
 r\log_2 B+\frac r2\log_2r
\]
が全 $i$ について成り立つ．従って全不変因子のビット長は
$O(r(\log B+\log r))$ 以下である．
\end{theorem}
\begin{proof}
$r\times r$ 小行列式全体の最大公約数を $\Delta_r$ とする．
左から整数の単模行列を掛けると，Cauchy--Binetの公式により
変換後の各 $r$ 次小行列式は変換前の小行列式の整数係数線形結合となる．
右からの単模行列についても同じであり，
逆変換も整数単模なので $\Delta_r$ は不変である．
Smith標準形では，階数 $r$ の非零小行列式は
$d_1\cdots d_r$ の一つだけだから $\Delta_r=d_1\cdots d_r$ である．

元の $A$ の任意の $r$ 次小行列 $C$ の各行のEuclidノルムは
$\sqrt r B$ 以下だからHadamardの不等式から
$|\det C|\le(\sqrt r B)^r$ である．
階数 $r$ だから非零の $r$ 次小行列式があり，
正の整数 $\Delta_r$ はその絶対値以下である．
よって $d_1\cdots d_r\le r^{r/2}B^r$．
各 $d_i$ は正整数なので $d_i\le d_1\cdots d_r$ となり，
対数を取れば表示式を得る．
\end{proof}
\begin{remark}
この評価は求める不変因子自身のビット長についての上界である．
中間計算で大きな係数を作らないSmith標準形算法の計算量まで
保証するものではない．
\end{remark}


\section{Smith標準形を使った完全な整数計算}
\begin{example}[行と列の可逆基本変形]
\[
 A=\begin{pmatrix}2&4\\6&8\end{pmatrix}
 \longrightarrow\begin{pmatrix}2&4\\0&-4\end{pmatrix}
 \longrightarrow\begin{pmatrix}2&0\\0&-4\end{pmatrix}
 \longrightarrow\begin{pmatrix}2&0\\0&4\end{pmatrix}.
\]
順に第二行から第一行の3倍を引き，
第二列から第一列の2倍を引き，
第二行の符号を変更した．従って
\[
 P=\begin{pmatrix}1&0\\3&-1\end{pmatrix},
 \quad Q=\begin{pmatrix}1&-2\\0&1\end{pmatrix},
 \quad PAQ=\operatorname{diag}(2,4).
\]
$\det P=-1$，$\det Q=1$ より両者は
$\mathrm{GL}_2(\mathbb Z)$ の元である．
しかも $2\mid4$ なのでSmith標準形である．
小行列式による検算では
$\Delta_1=\gcd(2,4,6,8)=2$，
$\Delta_2=|\det A|=8$ となり，
不変因子が $2,4$ と確かめられる．
\end{example}
\begin{proposition}[整数解の整除条件]
上の $A$ と $b=(b_1,b_2)^T$ に対し
$Ax=b$ が整数解を持つための必要十分条件は
$2\mid b_1$ かつ $4\mid(3b_1-b_2)$ である．
このとき $\mathbb Z^2/A\mathbb Z^2
\cong\mathbb Z/2\mathbb Z\oplus\mathbb Z/4\mathbb Z$ である．
\end{proposition}
\begin{proof}
$y=Q^{-1}x$ とすれば，$Ax=b$ は
$\operatorname{diag}(2,4)y=Pb$ に同値である．
$Pb=(b_1,3b_1-b_2)^T$ だから各対角成分での
整除性が必要十分である．$P,Q$ は整数格子の自己同型なので，
商加群の表示も対角行列の場合へ帰着する．
\end{proof}
\section*{到達確認問題と解答}
\noindent\textbf{問題．}
$Ax=(2,2)^T$ の整数解を全て求めよ．

\medskip\noindent\textbf{解答．}
$Pb=(2,4)^T$ なので $y=(1,1)^T$，
$x=Qy=(-1,1)^T$．$\det A=-8\ne0$ より解は一意である．
実際 $A(-1,1)^T=(2,2)^T$ となる．

\end{document}
